\subsubsection{Optimal Quoting and Inventory Control}

\begin{definition} \hlt{Inventory Imbalance}\\
Let $q_t$ be dealer inventory and $q_t^\star$ the target inventory. The imbalance is $q_t-q_t^\star$. Dealers quote to earn spread while keeping this imbalance financeable.
\end{definition}

\begin{table}[htbp]
\centering
\scriptsize
\stepcounter{table}
\renewcommand{\arraystretch}{1.13}
\setlength{\tabcolsep}{4pt}
\begin{tabularx}{\textwidth}{@{}>{\RaggedRight\arraybackslash}p{0.27\textwidth} >{\RaggedRight\arraybackslash}p{0.34\textwidth} >{\RaggedRight\arraybackslash}X@{}}
\hline
\textbf{Condition} & \textbf{Tactic} & \textbf{Purpose} \\
\hline
Inventories too low or & Raise bid price; increase bid size & Encourage clients to sell \\
clients are net buyers & Raise ask price; decrease ask size & Discourage clients from buying \\
& Take another trader's offer & Immediately raise inventories \\
& Buy a correlated instrument & Hedge inventory risk \\
\hline
Inventories too high or & Lower ask price; increase ask size & Encourage clients to buy \\
clients are net sellers & Lower bid price; decrease bid size & Discourage clients from selling \\
& Hit another trader's bid & Immediately lower inventories \\
& Sell a correlated instrument & Hedge inventory risk \\
\hline
\end{tabularx}

\vspace{0.25em}
{\footnotesize Table \thetable. Dealer tactics for managing inventories and order flow.}
\end{table}

\begin{principle} \hlt{Inventory Control Is Price Discovery}\\
A dealer's search for quotes that balance buy and sell flow also reveals market value. Fundamental value is often learned as a by-product of finding prices at which two-sided flow returns.
\end{principle}
